% GENERATED FILE. Edit canonical lessons, then run npm run generate:books.
% canonical-source-hash: fnv1a64:0141ba0c8b15112a
% canonical-lessons: MW-C20-hear-churma, MW-C20-churma, MW-C20-hear-chaay, MW-W20-ya, MW-C20-chaay, MW-R20-food-new-two, MW-C20-food-seven, MW-R20-script-close

\chapter{The Sweet and the Tea: One New Consonant}
\label{ch:mw-food-seven}
\hlchaptermodality{25}

\begin{chapteropening}
\textbf{By the end of this chapter:} \emph{I can add the sweet crumbled wheat dessert and tea to the food map after learning only \mw{य}.}

The last two source-attested food words begin by ear, add only \mw{य}, and complete a seven-word non-compensatory four-skill payoff.
\end{chapteropening}

\section[chūrmā]{Hear \emph{chūrmā} — sweet crumbled wheat}

\label{lesson:MW-C20-hear-churma}



\begin{quote}
Recall the five-word food payoff, say hard wheat rolls, then write wind.
\end{quote}

\subsection*{What to know first}
\begin{quote}
\emph{chūrmā} — \textbf{sweet crumbled wheat dessert}
\end{quote}

Three beats, the middle one long. Hear it twice and say it once. The source names it beside \emph{dāl} and \emph{bāṭī}; the three are often served together, and this lesson teaches the word rather than the custom.

\begin{tcolorbox}[breakable,colback=teal!4,colframe=teal!35!black,title={Before you move on}]
Source: \href{https://www.marwaripathshala.com/marwari-lesson-6-english}{Marwari Pathshala, Lesson 6}.
\end{tcolorbox}
\section[chūrmā]{\mw{चूरमा} — the sweet on the page}

\label{lesson:MW-C20-churma}



\begin{quote}
Recall the five-word transport payoff, then say the sweet crumbled wheat dessert and cloud.
\end{quote}

\subsection*{What to know first}
\begin{quote}
\textbf{\mw{चू}} + \textbf{\mw{र}} + \textbf{\mw{मा}} $\to$ \textbf{\mw{चूरमा}} — \emph{chūrmā} — sweet crumbled wheat dessert
\end{quote}

The \textbf{\mw{ू}} under \textbf{\mw{च}} is the long \emph{ū} met in \textbf{\mw{मिलसू}}, and \textbf{\mw{मा}} closes with the long \textbf{\mw{ा}} used since the first greeting.

\subsection*{Your turn}
Look, cover, wait five seconds, and write \textbf{\mw{चूरमा}}. Check that the \textbf{\mw{ू}} hangs below \textbf{\mw{च}} and not below \textbf{\mw{र}}.

\begin{tcolorbox}[breakable,colback=teal!4,colframe=teal!35!black,title={Before you move on}]
Source: \href{https://www.marwaripathshala.com/marwari-lesson-6-english}{Marwari Pathshala, Lesson 6}.
\end{tcolorbox}
\section[chāy]{Hear \emph{chāy} — tea}

\label{lesson:MW-C20-hear-chaay}



\begin{quote}
Say lentils, then write the sweet crumbled wheat dessert and cloud.
\end{quote}

\subsection*{What to know first}
\begin{quote}
\emph{chāy} — \textbf{tea}
\end{quote}

One long beat closing on a glide. It opens with the same consonant as \emph{chūrmā}. Hear it twice and say it once. The spelling waits one lesson, because its final consonant has not been taught yet.

\begin{tcolorbox}[breakable,colback=teal!4,colframe=teal!35!black,title={Before you move on}]
Source: \href{https://www.marwaripathshala.com/marwari-lesson-6-english}{Marwari Pathshala, Lesson 6}.
\end{tcolorbox}
\section[ya]{\mw{य} — the glide at the end of tea}

\label{lesson:MW-W20-ya}



\begin{quote}
Say tea and rain, then write hard wheat rolls.
\end{quote}

\subsection*{Script — one sign}
\begin{quote}
\textbf{\mw{य}} — \emph{ya}
\end{quote}

A left loop, a short arm, and a vertical stem on the right. Set it beside \textbf{\mw{भ}} of \textbf{\mw{भाई}}: \textbf{\mw{भ}} opens its left side into a hook that reaches the stem, while \textbf{\mw{य}} closes its left side into a loop.

\subsection*{Your turn}
Trace \textbf{\mw{य}} once. Copy it once. Write \textbf{\mw{भ}} beside it and say which one closes its left side. Cover both, wait five seconds, and write \textbf{\mw{य}} from the glide sound cue.

\begin{tcolorbox}[breakable,colback=teal!4,colframe=teal!35!black,title={Before you move on}]
Source: \href{https://www.unicode.org/charts/PDF/U0900.pdf}{Unicode Devanagari chart}.
\end{tcolorbox}
\section[chāy]{\mw{चाय} — tea on the page}

\label{lesson:MW-C20-chaay}



\begin{quote}
Say tea, then write \textbf{\mw{य}} and rain.
\end{quote}

\subsection*{What to know first}
\begin{quote}
\textbf{\mw{चा}} + \textbf{\mw{य}} $\to$ \textbf{\mw{चाय}} — \emph{chāy} — tea
\end{quote}

The \textbf{\mw{य}} carries no vowel mark of its own here and closes the word.

\subsection*{Your turn}
Look, cover, wait five seconds, and write \textbf{\mw{चाय}}. Check that the \textbf{\mw{ा}} sits after \textbf{\mw{च}} and that \textbf{\mw{य}} ends the word bare.

\begin{tcolorbox}[breakable,colback=teal!4,colframe=teal!35!black,title={Before you move on}]
Source: \href{https://www.marwaripathshala.com/marwari-lesson-6-english}{Marwari Pathshala, Lesson 6}.
\end{tcolorbox}
\section[Practice]{Retrieve the sweet and the tea}

\label{lesson:MW-R20-food-new-two}



\begin{quote}
Write \textbf{\mw{य}}, then say rain and write it.
\end{quote}

\subsection*{Your turn}
Hear both words in both orders, give each meaning, read two cards, then write both from sound.

\begin{tcolorbox}[breakable,colback=teal!4,colframe=teal!35!black,title={Before you move on}]
Source: \href{https://www.marwaripathshala.com/marwari-lesson-6-english}{Marwari Pathshala, Lesson 6}.
\end{tcolorbox}
\section[Practice]{Payoff — seven food words in four skills}

\label{lesson:MW-C20-food-seven}



\begin{quote}
Recall the earlier five-word food payoff and the three-word weather payoff.
\end{quote}

\subsection*{Your turn}
1. Identify seven heard words. 2. Produce seven words from meaning cues. 3. Match seven printed cards to meanings. 4. Write all seven heard words without a model.

Pass each skill separately. Seven food names are a vocabulary map; offering food, accepting it, asking for more, and refusing politely are all still untaught.

\begin{tcolorbox}[breakable,colback=teal!4,colframe=teal!35!black,title={Before you move on}]
Source: \href{https://www.marwaripathshala.com/marwari-lesson-6-english}{Marwari Pathshala, Lesson 6}.
\end{tcolorbox}
\section[Practice]{Close the writing loop for \mw{य}}

\label{lesson:MW-R20-script-close}



\begin{quote}
Recall the seven-word food payoff, then say vegetables and tea.
\end{quote}

\subsection*{Your turn}
Without a model, write \textbf{\mw{य}} and \textbf{\mw{ट}} from sound cues. Then write \textbf{\mw{चाय}}, \textbf{\mw{रोटी}}, \textbf{\mw{बाटी}}, \textbf{\mw{चूरमा}}, and \textbf{\mw{सबजी}} from dictation. Repair only the missed unit and repeat that word once.

\begin{tcolorbox}[breakable,colback=teal!4,colframe=teal!35!black,title={Before you move on}]
Sources: \href{https://www.marwaripathshala.com/marwari-lesson-6-english}{Marwari Pathshala, Lesson 6}; \href{https://www.unicode.org/charts/PDF/U0900.pdf}{Unicode Devanagari chart}.
\end{tcolorbox}
